\subsection{Limits of the relational fooling-set argument}\label{sec:natural-approaches}
% \subsection{Natural approaches that fail}\label{sec:natural-approaches}
A natural approach to lower bounds for full \CoR\ is to extend the earlier
exponential circuit lower bound~\citep[Theorems~6.1 and~6.12]{nakamura2022}
for the \defn{existential calculus of relations} $\Sigma_1\CoR$~\citep{nakamura2022,nakamura2023}, which is as expressive as existential \FOthree~\citep[Corollary~3.14]{nakamura2022}.
Here $\Sigma_1\CoR$ allows complement only on input relations and $\id$:
$\overline R;\one$ is permitted, but $\overline{R;\one}$ is not.
Although that lower bound was proved with formula-size games in the
style of Adler and Immerman~\citep{adler-immerman2003}, its proof can be
viewed as a relational fooling-set argument, which we recall first.

The \defn{fooling-set argument}~\citep{yao1979,lipton-sedgewick1981,kushilevitz-nisan1997,rao-yehudayoff2020}
is a standard technique for proving lower bounds in communication complexity.
For a Boolean function $f\colon A\times B\to\{0,1\}$, a
\defn{fooling set} is a set $F\subseteq f^{-1}(1)$ such that any two
distinct $(a,b),(a',b')\in F$ have a crossed input $(a,b')$ or
$(a',b)$ of value $0$. A \defn{1\nobreakdash-rectangle} is a product
$A'\times B'\subseteq f^{-1}(1)$. If a rectangle contains two distinct inputs of $F$,
it also contains their crossed inputs, one of which has value $0$; hence each
1\nobreakdash-rectangle contains at most one input of $F$, and covering
$f^{-1}(1)$ needs at least $|F|$ 1\nobreakdash-rectangles.

Consider Nakamura's query of size $O(k)$,
$\varphi_k(x,y)\coloneqq\exists z\,\bigwedge_{i=1}^k
(R_i(x,z)\lor R_i(z,y))$~\citep{nakamura2022}, where $k\ge2$.
%
For fixed $x$, $y$, and $z$,
let $I \coloneqq \set{i \in [k] : R_i(x,z)}$ and
$J \coloneqq \set{i \in [k] : R_i(z,y)}$.
Then $z$ is a witness precisely when $I \cup J = [k]$.
The Boolean function
$f_k(I,J)=1 \iff I\cup J=[k]$
has a fooling set
$F_k := \set{(I,I^c) : \varnothing \subsetneq I
\subsetneq [k]}$
of size $2^k-2$.

In~\citep[\S5]{nakamura2022}, each input $(I,I^c)$ is realized by the
three-element structure $\mathcal S_I$ consisting of a witness path
$u_I\xrightarrow{\,I\,}w_I\xrightarrow{\,I^c\,}v_I$ (the \emph{$I$-path}), so $\mathcal S_I\models\varphi_k(u_I,v_I)$.
A label $I$ means precisely the relations $R_i$ with $i\in I$ hold; unlisted entries, including loops, are false.

\begin{wrapfigure}{l}{0.3\linewidth}
  \centering
  \begin{tikzpicture}[
    >=Stealth,
    every node/.style={font=\small,inner sep=1pt}
  ]
    \node (u) at (0, 1.2) {$u^*$};
    \node (a) at (0, 0) {$a_{I,J}$};
    \node (c) at (1.5, 0.6) {$c_{I,J}$};
    \node (b) at (3.0, 1.2) {$b_{I,J}$};
    \node (v) at (3.0, 0) {$v^*$};

    \draw[->] (u) --
      node[above,sloped] {$I$} (c);
    \draw[->] (c) --
      node[above,sloped] {$I^c$} (b);
    \draw[->,dashed] (a) --
      node[below,sloped] {$J$} (c);
    \draw[->,dashed] (c) --
      node[below,sloped] {$J^c$} (v);
  \end{tikzpicture}
  \caption{\footnotesize $S_I\rtimes S_J$.
  Solid edges form $S_I$; dashed edges form $S_J$.}
  \label{fig:crossed-gadget}
\end{wrapfigure}
%
For each pair $(I,I^c),(J,J^c)\in F_k$ with $f_k(I,J^c)=0$,
let $S_I\rtimes S_J$ be the five-vertex structure in Figure~\ref{fig:crossed-gadget},
obtained by identifying the middle vertices of
copies of $S_I$ and $S_J$ as $c_{I,J}$,
with the other endpoints labeled as shown.
Let $\mathcal N$ be the union of these gadgets,
identifying all copies of $u^*$ and, separately,
all copies of $v^*$, while keeping all other
vertices distinct.
By construction, $S_I$ and $S_J$ embed into
$\mathcal N$, mapping $(u_I,w_I)$ to
$(u^*,c_{I,J})$ and $(w_J,v_J)$ to
$(c_{I,J},v^*)$, respectively.
Embedding is an injective map preserving and reflecting every input relation~\citep{hodges1993,marker2002}.
Yet every two-edge path from $u^*$ to $v^*$
is a crossed path with $f_k(I,J^c)=0$;
hence $\mathcal N\not\models\varphi_k(u^*,v^*)$.

$\Sigma_1\CoR$ terms are \emph{preserved under
embeddings}:
every term $V \in \Sigma_1\CoR$ and embedding
$\iota\colon\mathcal A\hookrightarrow\mathcal B$ satisfy the following:
\[
\tag{$\diamond$}\label{eq:mono}
\smash{\mathcal A\models V(u,v)\ \Longrightarrow\
\mathcal B\models V(\iota(u),\iota(v))}.
\]
%Thus, an accepted witness half remains accepted after embedding; this is exactly what makes relational crossing possible.
Now let $T$ be a $\Sigma_1\CoR$ term equivalent to $\varphi_k$.
For a composition $X;Y$ in $T$, let
$A_X\coloneqq\{I:\mathcal S_I\models X(u_I,w_I)\}$ and
$B_Y\coloneqq\{J^c:\mathcal S_J\models Y(w_J,v_J)\}$.
If $\mathcal N\not\models(X;Y)(u^*,v^*)$, then \textbf{$A_X\times B_Y$ is a 1\nobreakdash-rectangle}:
if some $(I,J^c)\in A_X\times B_Y$ had value $0$, then \eqref{eq:mono} applied to
the embeddings of $\mathcal S_I$ and $\mathcal S_J$ into $\mathcal N$ would give
$\mathcal N\models(X;Y)(u^*,v^*)$.
Hence $A_X\times B_Y$ contains at most one input of $F_k$.
For each $I$, $T$ contains a composition $X;Y$ with
$\mathcal S_I\models X(u_I,w_I)\land Y(w_I,v_I)$ and
$\mathcal N\not\models(X;Y)(u^*,v^*)$ (by induction on $T$, with converse pushed to the leaves).
Its rectangle contains $(I,I^c)$.
Thus $T$ has at least $|F_k|$ compositions.
This argument extends to circuits~\citep[\S6.3]{nakamura2022}.

\smallskip
\customparagB{What complement breaks.}
The fooling-set argument relies on the embedding-preservation property \eqref{eq:mono},
which fails in full \CoR\ due to unrestricted complement.
In fact, a simple \CoR\ term distinguishes  $\mathcal N$ from every $\mathcal S_I$.
Let $E\coloneqq\bigcup_{i=1}^k(R_i\cup R_i^\smile)$
and consider
$U\coloneqq\overline{(E\cap(E;\overline{\id}));\one}$.
The term $U$ holds at $(x,y)$ precisely when
$x$ has at most one $E$-neighbor.
Thus $U$ holds at $(u_I,v_I)$ in every
$\mathcal S_I$ but not at $(u^*,v^*)$ in
$\mathcal N$
because $u^*$ has distinct $E$-neighbors
$c_{I,J}, c_{I',J'}, \ldots$.


\if0
The argument relies on \eqref{eq:mono}, which fails in full \CoR\ already for the structures above.
Let $E\coloneqq\bigcup_{i=1}^k(R_i\cup R_i^\smile)$, and consider the term
$U\coloneqq\overline{(E\cap(E;\overline{\id}));\one}$, which has two compositions. The term $U$ holds at $(x,y)$ if and only if
$x$ has at most one $E$-neighbor. Hence $U$ holds at $(u_I,v_I)$ in every
$\mathcal S_I$ but not at $(u^*,v^*)$ in $\mathcal N$.
Thus these structures cannot certify an exponential lower bound for
\CoR, though $U$ does not define $\varphi_k$.
\fi
